% page 241 of the facsimile (image p-240 of the scan)
% block 167.6 x 233.3 mm ; left margin 34.4 mm
\pgc{12.700mm}{0.000mm}{
\l{\cel{55}233}
\l{}
\l{indolenta. Qua evitas penar. - DEFIS.}
\l{}
\l{indosar. (\sou{trans.})\cel{2}(\sou{kome}rco).Skribar e signatar la transfero-}
\l{\cel{3}formulo reverse di mandato, di kambio-letro. - DEFIS.}
\l{}
\l{induktar.\cel{2}(\sou{trans.}) I. (logiko). Rezonar, retroirante de la fakti}
\l{\cel{3}observita, a la lego (+leyo) quan oli konjektigas, de la efek-}
\l{\cel{3}ti a la kauzo. - II. (\sou{fiziko}). Produktar ta elektro-fluo qua}
\l{\cel{3}rezultas de la efiko da magneti o da konduktili trairata da}
\l{\cel{3}elektro-fluo (+korento). - DEFIRS.}
\l{}
\l{induktoro. (\sou{elektro}). (en\cel{1}\sou{mashino}\cel{1}elektrifala). La aparato qua}
\l{\cel{3}produktas la feldo induktera. - DEFIRS.}
\l{}
\l{indulgar. (\sou{trans., ulu, pri ulo}). Pardonar facile.sen postular}
\l{\cel{3}rigoroze exkuzi. - EFIS.}
\l{}
\l{induljenco. (\sou{teol.katol.}). Remiso, tota o parta, di la punisi}
\l{\cel{3}pro pekir, quan la ekleziokatolika grantas ula-kondicione}
\l{\cel{3}(fasto, prego, almono, e c.) - DEFIS.}
\l{}
\l{indulto.\cel{2}(\sou{yuro ekleziala}).- Povo, autoritato, grantita da la}
\l{\cel{3}papo, pri nominar ye ula benefici, o ye retenar ta benefici,}
\l{\cel{3}kontre la yuro komuna. - DEFIS.}
\l{}
\l{industrio. Ensemblo ek la arti, mestieri, qui transformas la}
\l{\cel{3}materii prima. - DEFIRS.}
\l{}
\l{\sou{indutar}.\cel{2}(\sou{trans.}) (ye\cel{1}\sou{ulo}).Kovrar ye farbo-strato. - F.}
\l{}
\l{inerta. I. Qua ne esas propre agiva. - II. (\sou{metaf}.) Sen-energia.}
\l{\cel{3}- EFIRS.}
\l{}
\l{\sou{infalibla}.\cel{2}Qua ne povas ne-eventor. - DEFIS.}
\l{}
\l{\sou{infama}. Stigmatizita da la opiniono sociala, da la lego. -DEFIS.}
\l{}
\l{\sou{infanto}. Homo qua evas inter 1 e 7 yari. - EFIS.}
\l{}
\l{\sou{\textquotedbl{}infante\textquotedbl{}}.\cel{2}En Hispania ed en Portugal, nomo di ta ek la princi}
\l{\cel{3}qua naskis nemediate pos la senioro.}
\l{}
\l{\sou{infantrio}.\cel{2}(\sou{militarto}). Ensemblo ek armeani qui marchas, e}
\l{\cel{3}kombatas\cel{2}pedo-stacante. - DEFIRS.}
\l{}
\l{\sou{infarkto}. (\sou{patol.}) Infiltreso da tisuo ye ekfluo da sango. - DEFIS}
\l{}
\l{infektar.\cel{2}(\sou{trans.})\cel{2}Impregnar ye jermi nociva. - DEFIRS.}
\l{}
\l{\sou{infera}r.\cel{2}(\sou{trans., de, ek)}\cel{3}Ektirar (konsequajo). - EFIS.}
\l{}
\l{\sou{inferi}ora. (en\cel{1}\sou{hierarkio}). Qua okupas rango suba, infra.- DEFIS.}
}
